\subsection{Religion und Religiosität}    \label{subsec_2.1.1}

Die Religion verleiht vielen Menschen ihrem Leben Sinn und Zweck. Selbst in stark säkularisierten Ländern wie Westeuropa, wie vom \textcite{PRC} berichtet, gaben im Median 39\% an, dass die Religion ihnen dieses Gefühl vermittelt. Zudem bietet sie Trost in Zeiten der Trauer und moralische Orientierung bei schwierigen Entscheidungen. Viele Menschen sehen das Hauptziel von organisierter Religion darin, Menschen zu vereinen und zusammenzubringen \parencite{PRC}. 

Unter organisierter Religion fallen Religionsgemeinschaften, die ein einheitliches und förmlich festgelegtes Glaubenssystem und Rituale besitzen. Diese Systeme und Rituale sind oftmals durch eine offizielle Doktin, eine hierarchische Führungskraft und einer Kodifizierung der richtigen und falschen Verhaltensweisen festgelegt. Diese unterscheiden sich von nichtorganisierter Religion. Die meisten Personen nichtorganisierter Religion glauben an Gott oder eine höhere Macht, sie haben aber keine Beziehung zu einer organisierten Religionen und umfassen moderne Spiritualitätsformen \parencite{OrgaRel}. Davon zu unterscheiden ist die organisierte und nichtorganisierte Religiosität. Religion bezieht sich auf die äußerliche Beschreibungsebene, wobei Religiosität das innere Erleben beschreibt \parencite{ReligionReligiosität}. Unter organisierter Religiosität werden solche Praktiken verstandne, die in einem sozialen Umfeld gelebt werden. Beispiel dafür sind der Besuch von Gottesdiensten, die Teilnahme an religiösen Gruppenritualen, gemeinsames Bibelstudium und mehr. Das private Gebet Zuhause, Meditation, das Hören religiöser Musik sind hingegen Beispiele für nichtorganisierte Religiosität \parencite{ReliSpiriModell, DUREL_Val}. Im Rahmen dieser Studie wurde zudem auch die intrinsische und extrinsische Religiosität beobachtet. Eine Person mit hoher intrinsicher Religiosität lebt ihre Religion und übt diese aus innerer Überzeugung heraus aus. Sie verfolgt ihre Religion als ultimatives Ziel an sich, unabhängig von äußeren Belohnungen oder Zielen. Das Ziel ist die spirituelle Verbundenheit und das persönliche Wachstum. Extrinsische Religiosität hingegen sieht diese nur als Mittel zum Erreichen wichtigerer Ziel, wie finanzieller Erfolg, sozialer Satus oder Komfort. Eine Person mit hoher extrinsischer Religiosität stellt diese hauptsächtlich nur zur Schau, um irdische Vorteile zu genießen \parencite{DUREL_Val}. 

Eine andere Studie des \textcite{PRC2} untersuchte, ob Religion eine positive, negative oder gemischte Assoziation mit verschiedenen Indikatoren des individuellen und gesellschaftlichen Wohlbefinden aufweist. Sie fanden, dass in den USA und weiteren Ländern die regelmäßige Teilnahme an einer Religionsgemeinschaft mit einem höheren Maß an Zufriedenheit verbunden ist. Das Institut postulierte die Hypothese, dass Gesellschaften, in denen das religiöse Engagement abnimmt, möglicherweise mit einem Rückgang des individuellen und gesellschaftlichen Wohlbefindens konfrontiert sind. Die Religionszugehörigkeit scheint die Wahrscheinlichkeit für persönliches Glück und Wohlbefinden nicht zu beeinflussen \parencite{PRC2}.

Schon lange wird in der Wissenschaft diskutiert, wie Religion zu definieren ist. Einige stellen die Frage, ob der Begriff überhaupt eine stabile Kategorie menschlicher Erfahrung beschreibt \parencite{Taxonimie, ReliBegriff}. \textcite{BritannicaDef} gibt zwei mögliche Definitionen für Religion: Religion ist \enquote {der Glaube an einen Gott oder an eine Gruppe von Göttern}, Religion als \enquote {ein organisiertes System von Glaubensvorstellungen, Zeremonien und Regeln, die der Verehrung eines Gottes oder einer Gruppe von Göttern dienen}. Es gibt aber noch viele weitere Definitionen, die sich auf unterschiedliche Aspekte fokussieren, wie es bei \textcite{OxfordDic} zu finden ist. Selbst die Etymologie des Wortes hat mehr als einen möglichen Ursprung. In jenen sieht \textcite{OxfordDic} dennoch eine gemeinsame Richtung. Alle Definitionen über die Jahrungerte hinweg sowie die Etymologie und deren Bedeutungen lenken die Aufmerksamkeit auf eines der wichtigsten Merkmale der Religion. Sie ermöglicht es den Menschen, die Welt zu begreifen und sich im Leben zu orientieren \parencite{Reli}.
Aus der Sicht der Soziologie haben Religionen schon immer einen Wert gehabt, primär weil sie dem Überleben und der Selektion dienen. Die Religion ist ein hoch organisiertes Schutzsystem. Die Evolutionspsychologie zeigt, dass Religion und Spiritualität möglicherweise evolutionäre Vorteile bieten, indem sie das Überleben durch soziale Bindungen und Vertrauen fördern. Rituale und religiöse Praktiken können als aufrichtige Signale dienen, die die Kooperation und den Zusammenhalt innerhalb von Gemeinschaften stärken \parencite{NorenzayanShariff2008}. Dies trägt zur Erhöhung des individuellen Wohlbefindens und der sozialen Stabilität bei. Für Jahrtausende haben Menschen ein gewisses Religionssystem gehabt, worin die Menschen ihr Leben erfolgreich gelebt haben. Bereits vor 12.000 Jahren bauten frühe Jäger und Sammler beeindruckende Monumente wie Göbekli Tepe, was auf eine weit zurückreichende Glaubenskultur hinweist \parencite{Joseph2011}. Diese historischen Entwicklungen zeigen, dass spirituelle Rituale und Überzeugungen tief in der menschlichen Geschichte verankert sind. Heutzutage zählen sich mehr als drei Viertel der Weltbevölkerung einer Religion zugehörig, unabhängig wie viel oder wenig sie damit zu tun haben \parencite{OxfordDic}.

Spiritualität ist ein sehr komplexes Konstrukt, das sich durch eine wahrgenommene Verbundenheit zum Heiligen oder Transzendenten auszeichnet. Sie geht über das hinaus, was wir sehen oder verstehen können, und beinhaltet existenzielle Fragen wie den Sinn des Lebens \parencite{Spiritualität}. Religion hingegen ist ein System von spezifischen Anschauungen und Praktiken, die sich auf heilige Dinge beziehen und oft institutionell organisiert sind \parencite{Durkeim1981}. Die Verbindung von Spiritualität und Religion zeigt sich darin, dass Spiritualität sowohl innerhalb von Religion existieren kann als auch außerhalb von ihr \parencite{Büssing2008}.

Wie ist die Religion aber von der Spiritualität zu unterscheiden? Es gibt unterschiedliche Modelle, die versuchen zu erklären, wie Religiosität und Spiritualität miteinander interagieren. \textcite{Spiritualität} definiert Spiritualität als eng verbunden mit dem Übernatürlichen und Religiösen. Nach seiner Definition ist die Suche nach dem Transzendenten teil von Spiritualität. Ein Nichtgläubiger gelangt durch sie zum Hinterfragen, dann zum Glauben und schließlich zur Hingabe. Diese Definition von Spiritualität findet sich im Modell des traditionellen Verständnisses von Spiritualität von \textcite{ReliSpiriModell}. In ihm ist die Spiritualität ein Teil von Religion und Religiosität. \textit{Spirituel} sein dient in diesem Modell zur Beschreibung einer tief religiösen Person. Beim moderneren Verständnis hat sich die bedeutung von Spiritualität erweitert und umfasst diejenigen, die sich als religiös bezeichenen und jene, die es nicht tun, sich aber als spirituel empfinden. Diese Veränderung hat zu vielen Problemen geführt, um Spiritualität gut definieren zu können, sie von Religiosität soweit zu trennen, dass Personen, die sich als spirituel, aber nicht religiös identifizieren nicht vernachlässigt werden und letzendlich Spiritualität zu operationalisieren.

Die erfolgreich etablierten Schutzsysteme geben den Menschen Gewissheit und Sicherheit auch andere Dinge zu verfolgen. Dadurch haben sie die Möglichkeiten ihre eigene Natur und Gesellschaften, wie die Welt, die sie umgibt zu erforschen. Die Untersuchung der eigenen Natur schaut zu was der Körper fähig ist zu erleben. Worin liegt sein Potential? Einige Religionen legen den Schwerpunkt aufs Innere des Menschen. Der Buddhismus beispielsweise sucht die Wahrheit im Körper und im Sinne der Erleuchtung, Friede und Leere. Andere Religionen legen den Schwerpunkt aufs Außen. Welchen Wert haben die Dinge außerhalb des Körpers und die Beziehungen mit anderen; deren absoluter Höhepunkt die Gemeinschaft und Vereinigung mit Gott ist? Aus Systemen, die den Wert von Beziehungen untersuchen sind das Judentum, das Christentum und der Islam entsprungen. Dennoch ist festzuhalten, dass beide Systemarten auch den Schwertpunkt des anderen mit betrachtet. Beide erkennen auch die Realität des Bösen undes Schlechten und alle wissen, dass die Erforschung durch den Menschen allein nicht ausreicht und das göttlich dabei unterstützt und fördert \parencite{OxfordDic}. Obwohl verschiedene Religionen wahrhaftig unterschiedlich sind, schützen sie Menschen und vermitteln ihnen wichtige Erkenntisse in die Möglichkeiten des menschlichen Lebens. All das zeigt, dass Religionen Schutzsysteme sind, die essenzielle Aspekte des menschlichen Lebens und Wohlbefindens bewahren. Die Art und Weise, wie sie funktionieren und wirken, sind eng mit dem Konzept under Praxis der Religion verbunden.

\subsubsection{Auswirkung auf die Gesundheit}    \label{subsubsec_2.1.1.1}
Religionen basieren auf einer eigenen, spezifischen religiösen Logik und beschäftigen sich mit grundlegenden, aber nicht abschließend entscheidbaren Fragen. Diese umfassen Themen wie die Stellung des Menschen in der Welt, das Leben nach dem Tod und die Definition eines „guten Lebens“ sowie moralische Pflichten gemäß den Werten der spezifischen Religionsgemeinschaft. Aufgrund dessen haben Religionen einen erheblichen Einfluss auf Erziehung, Bildung, Handeln und auch auf die Gesundheit und gesundheitsrelevantes Verhalten. Sie bieten ihren Mitgliedern eine umfassende symbolische Weltanschauung und eine Vielzahl spezifischer religiöser Praktiken, die darauf abzielen, den Menschen Orientierung, Sinn und Halt im Leben zu geben. Dieser ganzheitliche Anspruch spiegelt sich in der Schaffung und Weitergabe eines reichen symbolischen Universums wider, das oft als eine lebensprägende „Erzählung“ oder Narration fungiert. Diese symbolischen und narrativen Strukturen haben häufig eine direkte Relevanz für den Alltag und beeinflussen das Verhalten der Gläubigen. Dazu gehören Vorstellungen von der richtigen sozialen Ordnung, Anstand und Moral, Machtverhältnissen und Ohnmacht, sowie Selbst- und Weltbilder. Zudem prägen Religionen auch körper- und geschlechtsbezogene Rollenbilder sowie Ansichten zu Ernährung und Sexualität \parencite{Hemel2023}. Die Beziehung zwischen Religiosität und Gesundheit ist jedoch komplex und ambivalent. Es gibt sowohl positive als auch negative Auswirkungen, die in der wissenschaftlichen Literatur diskutiert werden. Unumstritten ist, dass schon bevor der Mensch Psychosomatik als solche definierte, in religiösen Weltsichten die körperliche und seelische Gesundheit als verbunden betrachtet wurde und wird.

Sowohl \textcite{Hemel2023} als auch \textcite{Jakobs_2018} betonen, dass Religiosität positive Effekte wie die Förderung sozialer Unterstützung und Gemeinschaftsgefühls sowie die Sinnstiftung hat. Sie fördert die soziale Unterstützung, welche wiederum positive Effekte auf die Gesundheit hat. Außerdem verstärken Religiosität die emotionale und behaviorale Selbstregulation, was zu einem besseren Gesundheitsverhalten und zu günstigeren physiologischen Prozessen führt. Diese Effekte tragen dazu bei, die Häufigkeit von Erkrankungen zu reduzieren und die Mortalität zu senken \parencite{Krause2015}. Religiöse Überzeugungen können Menschen helfen, Krisen zu bewältigen und dem Leben einen tieferen Sinn zu geben. Andererseits können strikte religiöse Normen zu Schuldgefühlen und Angst führen, und extreme religiöse Praktiken können soziale Isolation und psychische Belastungen verursachen.

\textcite{Hemel2023} hebt hervor, dass soziokulturelle Praktiken, die durch Religion vermittelt werden, Schutz und Unterstützung bieten können. Dieser Effekt sei insbesondere bei Frauen gegeben. Allerdings betont er auch, dass restriktive Normen und Rollenbilder die Autonomie und das Wohlbefinden, insbesondere von Frauen, beeinträchtigen können.

\textcite{Jakobs_2018} konzentriert sich auf die individuelle religiöse Erfahrung und deren Bedeutung für die allgemeine psychische Gesundheit. Er unterstreicht, dass eine positive Religionspädagogik die positiven Effekte der Religiosität verstärken kann. Zusätzliche spezifische negative Effekte gegenüber den allgemeinen Aussagen führt Jakobs nicht an.

\textcite{Krause2015} beschreibt, dass das Gebet als häufige Form der spirituellen Praxis sowohl positive als auch negative Auswirkungen auf die psychische Gesundheit haben kann. Positive Auswirkungen umfassen die Förderung von Wohlbefinden und die Reduktion von Stress. Negative Auswirkungen können auftreten, wenn Gebete von Angst und Hader geprägt sind oder wenn sie als einzige Bewältigungsstrategie verwendet werden, ohne andere unterstützende Maßnahmen zu ergreifen. Religiosität wirkt direkt und in ihrer Wirkung positiv auf wichtige Faktoren wie gesundheitsbezogene Lebensqualität, mentale Gesundheit, körperliche Erkrankungen und eine geringere Sterberate aus. Studien zeigen, dass religiöse Praktiken das Auftreten von Herzerkrankungen, Bluthochdruck, Infarkten und Schmerzen verringern können. Außerdem fördern sie positives Gesundheitsverhalten und das Langzeitüberleben. Hoffnung ist ein wichtiger Faktor für das psychische Wohlbefinden und wirkt als Puffer gegen psychische Erkrankungen. Sie hängt mit einem höheren Selbstwertgefühl, mehr positiven Gefühlen und einem besseren subjektiven Wohlbefinden zusammen. In Zeiten von Krankheit und Krisen stellt Hoffnung eine wertvolle Ressource dar \parencite{Krause2015}. 







%\subsubsection{Konfessionen}    \label{subsubsec_2.1.1.2}


%Der Mensch hat schon immer nach einer höheren Macht gesucht und sie Götter genannt. Sie wurden als allmächtige und allwissende Wesen verehrt. Jede Kultur hatte ihre individuellen Gottheiten, die ihre lokale Religion bildeten. Heutzutage werden die meisten, wie beispielsweise die Götter des antiken Griechenlands davon nicht mehr verehrt. Heutzutage bekennen sich die meisten Menschen zu einer der fünf großen Religionen: das Judentum, das Christentum, der Islam, der Hinduismus und der Buddhismus. Diese Konfession unterschiedlichen sich untereinander. Sie haben unterschiedliche Glaubensgrundsätze. In diesem Unterkapitel werden die fünf weltreligionen skiziert um einen groben Überblick zu schaffen \parencite{ReliARD}.

%Das \textit{Judentum} ist die älteste monotheistische Weltreligion und ist über 3.000 Jahre alt. Ihren Gott nennen sie auch Jahwe. Die Juden glauben, dass ihr Gott einen Erlöser schicken wird, um den Frieden zu bringen. Des Weiteren glauben sie an ein Leben nach dem Tod. Ursprünglich lebten die Juden in dem heutigen Isreal. Durch immer wieder aufkehrende Verfolgungen verließen sie ihr Land und verteilten sich somit über die ganze Welt. Heutzutage gibt es ca. 14 Millionen Juden auf der Welt. Israel und dessen Hauptstadt Jerusalem ist für sie das Heilige Land. Der Davidstern bestehend aus zwei Dreiecken ist das Symbol der Juden und repräsentiert die enge Verbindung zwischen Gott und den Menschen. Die Synagoge ist das Gotteshaus der Juden. An jedem Sabbat, der Freitagabend beginnt und bis Samstagabend andauert, sollen sich die Juden ausruhen, beten und aus der Tora lesen, das in der Synagoge traditionell auf Hebräisch vorgelesen wird. Die Tora, bestehend aus den fünf Büchern Moses und ist der erste Teil ihres heiligen Buches, der Tanach. Die anderen beiden Teile sind die Bücher der Propheten Newiin und die Schriften Ketuvims. Neber dem Sabbat feiern die Juden acht Tage lang das Passah-Fest im Frühling. Es gedenkt der Befreiung aus Ägypten. Ein Erkennungsmerkmal vieler jüdischer Jungen und Männer ist die Kippa. Sie symbolisiert die Ehrfurcht vor Gott \parencite{ReliFocus, ReliARD}.
%Zusammenhalt in Gemeinde und Familie sehr wichtig

%Aus dem Judentum entstand das \textit{Christentum}. Jesus Christus wird als die zweite Person der Dreifaltigkeit gesehen. Die erste Person ist Gott$-$Vater, gefolgt von Gott$-$Sohn und vom Gott$-$Heiligen Geist. Jesus ist Mensch geworden und lebte vor ca. 2.000 Jahren unter den Menschen. Er war Jude. Seine Taten und Worte sind im Neuen Testament in der Bibel dokumentiert. Die Bibel ist das Heilige Buch der Christen und besteht aus 66 bis 73 Büchern. Das Alte Testament erzählt unter anderem über die Schöpfung der Welt und die Anfänge der Menschheit. Neber den Worten und Taten Jesu bis hin zur Auferstehung umfasst das Neue Testament auch Informationen über die Zeit danach. Für die Anhänger Jesus, war er der Erlöser, auf den die Juden gewartet haben und sie schlossen mit Gott einen neuen Bund. Heutzutage bildet das Christentum die größte Weltreligion mit mehr als 2 Milliarden Bekennern. Sie ist in vier Hauptgruppen unterteilt: die römisch-katholische, die protestantische, die orthodoxe und die anglikanische Kirche. Für alle ist das Symbol des Kreuzes eine Erinnerung an Jesu Tod und seine Auferstehung nach drei Tagen. Auch die Feste im Christentum beziehen sich auf Jesu Leben wie beispielsweise Weihnachten (Jesu Geburt), Karfreitag (Leid und Tod Jesu) und Ostern (Auferstehung Jesu). Für viele Christen ist der Besuch der Kirche, dem Gotteshaus der Christen, an den Feiertagen verpflichtend. Im Christentum ist Ostern das wichtigste Fest, weil Jesus mit seiner Auferstehung den Tod bezwungen hat. Für alle Christen ist die Taufe sehr wichtig, denn diese nimmt die Person in die christliche Gemeinde auf. Bei den meisten Christen wird dies Zusprache zu Gott mit der Konfirmation bei den Protestanten und der Ersten Heiligen Kommunion und später der Firmung bei den Katholiken erneuert und bestätigt. Die katholische Kirche besitzt des Weiteren über ein Oberhaupt, den Papst. Ähnlich wie die Juden, glauben auch die Christen an ein Leben nach dem Tod. Christen glauben an die Auferstehung der Toten und das ewige Leben \parencite{ReliFocus, ReliARD}.

%Die jüngste, aber mit 1,6 Milliarden Anhängern zweitgrößte der fünf Weltreligion ist der \textit{Islam}, der vor ca. 1440 Jahre vom Propheten Mohammed aus dem Judentum und Christentum heraus erschaffen wurde. Die Muslimer nennen ihren Gott Allah. Es existieren zwei Hauptgruppen von Muslimen, die Sunniten und die Schiiten. Der Koran, aufgeteilt in 114 Kapitel auch Suren genannt, das heilige Buch der Muslimer, beinhaltet Allah's Botschaft und alle Regeln, die ein Muslimer befolgen muss. Das ist auch das Ziel der Muslimer, vollkommen nach diesen Gesetzen leben, wie es ihr Gott erwartet. Muslimer glauben, dass Gott den Menschen Propheten geschickt hat. Neben Noah, Abraham, Mose und Jesus, die auch für Christen und teilweise auch für die Juden wichtige Personen sind, ist für die Muslimer auch Mohammed einer der wichtigsten Propheten. Sein Geburtstag ist zusammen mit dem Ende des Ramadan und dem Opferfest eines der wichtigsten Feste im Islam. Dadurch, das Mohammed in Mekka und Medina geboren und aufgewachsen ist, sind dies wichtige Städte in der islamischen Religion. In Mekka steht das wichtigste Heiligtum der Muslime, die Kaaba, ein großer schwarzer Würfel, in dessen Richtung Muslimer fünf Mal täglich zu festgelegten Zeiten beten. Das Gotteshaus im Islam heißt Moschee und ihr Symbol ist der Halbmond, der an die Eroberung Istambuls durch muslimische Krieger erinnern soll, die sich das Symbol damals aneigneten. So wie die Juden und Christen glauben auch die Muslima an das ewige Leben nach dem Tod. Sie glauben, dass Allah entscheidet, ob sie in die Hölle oder in das Paradies gelangen, abhängig davon, wie sie ihr Leben gelebt haben \parencite{ReliFocus, ReliARD}. 

%Die älteste der fünf Weltreligionen ist der \textit{Hinduismus}, der ca. 4.000 v.Ch. entstand. Konträr zu den bis jetzt genannten monotheistischen Religionen ist der Hinduismus eine polytheistische Religion. Im Hinduismus existieren Millionen verschiedene Götter mit ihren unterschiedlichen Eigenschaften. Die meisten der ca. 800 Millionen Hindus verehren nicht alle existieren Gottheiten an, sondern tuhen dies nur für persönlich ausgewählte Götter. Einige Hindus glauben auch, dass es nur ein Gott ist, der sich in vielen zeigt. Die meisten Hindus leben in Indien. Dort gibt es auch die meisten Tempel, die Gotteshäuser im Hinduismus, die meist nur einer Gottheit geweiht sind. Im Gegensatz zum Judentum, Christentum und dem Islam gibt es beim Hinduismus nicht ein spezifisches heiliges Buch. Für die Hindus sind die Veden sehr wichtig. Sie sind eine Sammlung von Liedern, Gedichten und Geschichten. Wenn aus den heiligen Schriften gelesen wird, beginnt jeder Vers mit der Silbe OM. Sie symbolisiert das göttliche Prinzip und soll dem Leser helfen eine Verbindung zu Gott herzustellen. Diese Silbe ist auch das Symbol der Hindus. Ein weiterer Unterschied zu den genannten monotheistischen Religionen ist der Glaube der Hindus an die Wiedergeburt. Für sie ist die Seele zwar auch unsterblich, sie glauben aber, dass sie nach dem Tod in eine andere körperliche Gestalt übergeht. Die neue Gestalt hängt davon ab, wie die Person in ihrem früheren Leben gelebt hat. Hindus glauben, dass eine Seele in einem Menschen, einem Tier oder in einer Gottheit wiedergeboren werden kann. Aus diesem Grund sind viele Hindus vegetarier. Jene, die nicht auf Fleisch verzichten, essen jedoch keine Rindfleisch, denn Kühe gelten im Hinduismus als heilige Tiere und es ist jedem Hindu verboten diese zu essen. Diesen Kreislauf des Sterbens und Wiedergeboren werden, versuchen Hindus zu entkommen, durch Erleuchtung. Wenn sie dies schaffen, dann gelangt ihre Seele ins Nirwana. Wenn Hindus sterben führen die Angehörigen ein traditionelles Beerdigungsritual durch, das mehrere Tage andauert. Neber diesem wichtigen Ritual gibt es auch viele wichtige Feste im Hinduismus, wie das Holi-Fest, oder auch Frühlingsfest genannt, und das Kumbh Mela, oder Krugfest auf deutsch \parencite{ReliFocus, ReliARD}.

%Auch im \textit{Buddhismus} sthet wie im Hinduismus die Erleuchtung im Mittelpunkt. Gegründet wurde diese Religion vor ca. 2.500 Jahren durch Buddha, der eigentlich Siddhartha Gautama hieß. Im Hinduismus gilt er als erleuchtet. Das Dharma$-$Rad, das Symbol der Buddhisten soll seine Lehren darstellen. Des Weiteren symbolisiert es den ewigen Kreislauf des Lebens, Samsara, das kein Anfang und kein Ende hat. Die ca. 360 Millionen Buddhisten glauben, dass sie nach dem Tod irgendwann wiedergeboren werden. Diejenigen, die Erleuchtung erlangen, können diesen ewigen Kreislauf entweichen und ins Nirwana gelangen, einem heiligen Seelenzustand. Im Buddhismus wird geglaubt, dass wer Buddha's Lehre folgt, das Leid des eigenen Lebens beendet. Einer dieser Lehren ist, dass Buddhisten logisch und vernünftig handeln sollen; Dinge zu hinterfragen, statt sie blind anzunehmen. Diese und weitere Lehren sind in dem heiligen Buch der Buddhisten, dem Pali$-$Kanon, hinterlegt. Der Kanon besteht aus drei Teilen. Der erste Teil umfasst die Lebensregeln für buddhistische Mönche und Nonnen. Im zweite Teil sind die Reden und weitere Geschichten über Buddha dokumentiert, dessen Lehren im dritten Teil erklärt werden. Auch die Feste sind auf Buddah ausgerichtet und erinnern an sein Leben. Das wichtigste Fest, Vesakh, feiert seine Gerburt. Für ihre Gebete und Meditation gehen Buddhisten in die reich verzierten Tempel oder in Gebetshallen. Es gibt auch kleinere Gebäude mit überresten von Buddah oder anderen wichtigen Buddhisten, so genannte Stupas, in denen Buddhisten beten und meditieren können \parencite{ReliFocus, ReliARD}.